\subsection{Unit-connector ordinary eighth moments for even terminals of minimum one}\label{proof:Q5-C061}

\paragraph{Q61.0 The terminal and its directions.}
Use the fixed finite model of Appendix~\ref{proof:Q5-C064}, with
\(Jv=\allowbreak (\E v)\one\), \(A=\allowbreak J+K\), \(T=\allowbreak J-K=\allowbreak A^*\), \(K^*=\allowbreak -K\),
\(|K(x,\allowbreak y)|\le\allowbreak 1\), \(g=\allowbreak K\one\), \(p=\allowbreak 1+g\), \(q=\allowbreak 1-g\),
\(d=\allowbreak \E g^2\), \(\delta_j=\allowbreak \|K^j\one\|_2^2\),
\(\theta=\allowbreak 13/32\), \(\kappa=\allowbreak 19/32\), \(f_n=\allowbreak A^n\one\),
\(h_n=\allowbreak Tf_{n-1}\), and \(h_1=\allowbreak q\).
After maximal equal pairing, suppose
\(1=\allowbreak a_1<\allowbreak b_1<\allowbreak a_2<\allowbreak b_2<\allowbreak \cdots<\allowbreak a_k<\allowbreak b_k\), \(k\ge\allowbreak 2\).
For the genuine equal-pair product \(E_0=\allowbreak \prod_\ell f_\ell h_\ell\le\allowbreak 1\), set
\begin{equation}\tag{Q61.0.1}\label{eq:Q61:0.1}
\boxed{F=E_0qf_{b_1}\prod_{i=2}^kh_{a_i}f_{b_i},\qquad H=pTF.}
\end{equation}
Each \(f_n\)-arm points outward; each \(h_n\)-arm has only its first
edge inward. The terminal unit arm points inward. A single connector
edge points from the terminal to the new root, whose unit arm points
outward. This choice depends only on the arm list. For all finite
integer lengths, \(k\ge\allowbreak 2\), legal kernels and equal-pair multiplicities,
\begin{equation}\tag{Q61.0.2}\label{eq:Q61:0.2}
\boxed{\E[(pTF)^8]\le e^{-63d/50}\le1-d/2.}
\end{equation}
The proof includes \((1,\allowbreak 2,\allowbreak a,\allowbreak c)\) with \(a\ge\allowbreak 4\), \((1,\allowbreak 3,\allowbreak a,\allowbreak c)\) with
\(a\ge\allowbreak 4\), every \(b_1\ge\allowbreak 4\), and \((1,\allowbreak 2,\allowbreak 3,\allowbreak c)\), in the ordered
direction \eqref{eq:Q61:0.1}.

\paragraph{Q61.1 Finite identities and boundary cases.}
After deletion of zero-mass atoms, skewness gives
\(|g(x)|\le\allowbreak 1-\mu(x)\) and \(p(x),\allowbreak q(x)\ge\allowbreak \mu(x)>\allowbreak 0\).
The finite calculus gives
\begin{equation}\tag{Q61.1.1}\label{eq:Q61:1.1}
0\le d\le1/3,\quad\|K\|_{2\to2}^2\le\theta,\quad
\delta_{j+1}\le\theta\delta_j,\quad |Kv(x)|\le\|v\|_2,
\end{equation}
and
\begin{equation}\tag{Q61.1.2}\label{eq:Q61:1.2}
\E K^{2j}\one=(-1)^j\delta_j,\quad\E K^{2j+1}\one=0,\quad
0\le f_n,h_n\le2,\quad f_nh_n\le1.
\end{equation}
These are identities in \(L^2(\mu)\). When \(d=\allowbreak 0\), \(g=\allowbreak 0\) and
\(A\one=\allowbreak T\one=\allowbreak \one\); integration from the leaves gives every tree
message equal to one even if \(K\ne0\). Below take \(d>\allowbreak 0\) where needed.

\paragraph{Q61.2 The \(K^2\) row input.}
The finite three-label and random-sign argument in
\eqref{eq:Q49:2.1}--\eqref{eq:Q49:2.3} proves
\fitdisplay{
\|Kv\|_2^2\le1/3\quad(|v|\le1),\qquad
|K^2w(x)|^2\le\|w\|_2^2/3\quad(w\text{ real}).
}{}
The first estimate is from \(L^\infty\) to \(L^2\), and the second
is a row estimate for \(K^2\). The spectral estimate used here is
\(\|K\|^2\le\allowbreak \theta\).

\paragraph{Q61.3.1 A centered variable in a bounded interval.}
Let a finite real variable \(Y\) have \(\E Y=\allowbreak 0\) and take values in
an interval of length \(\ell\). For \(\psi(s)=\allowbreak \log\E e^{sY}\), the
exponentially tilted law has the same support, and
\(\psi''(s)=\allowbreak \operatorname{Var}_s(Y)\le\allowbreak \ell^2/4\).
Indeed, if \(c\) is the interval midpoint, then
\(\operatorname{Var}_s(Y)\le\allowbreak \E_s(Y-c)^2\le\allowbreak \ell^2/4\).
Since \(\psi(0)=\allowbreak \psi'(0)=\allowbreak 0\), two integrations give, for every \(s\ge\allowbreak 0\),
\begin{equation}\tag{Q61.3.1}\label{eq:Q61:3.1}
\E e^{sY}\le e^{s^2\ell^2/8}.
\end{equation}

\paragraph{Q61.3.2 Hermite comparison with either sign of the mean.}
Suppose the same finite real variable \(W\) satisfies
\(|W|\le\allowbreak  U\), \(\E W^2\le\allowbreak  U^2/2\), \(U\ge\allowbreak 0\).
For \(s\ge\allowbreak 0\), put \(m=\allowbreak \E W\), \(m_+=\allowbreak \max(m,\allowbreak 0)\). We prove
\begin{equation}\tag{Q61.3.2}\label{eq:Q61:3.2}
\boxed{\E e^{sW}\le\exp\{sm_++9s^2U^2/32\}.}
\end{equation}
The cases \(U=\allowbreak 0\) or \(s=\allowbreak 0\) are immediate. For \(U,\allowbreak s>\allowbreak 0\), let
\(f(w)=e^{sw}\). Its Hermite majorant is
\begin{equation}\tag{Q61.3.3}\label{eq:Q61:3.3}
Q(w)=f(a)+f'(a)(w-a)+B(w-a)^2,\qquad
B=\frac{f(U)-f(a)-f'(a)(U-a)}{(U-a)^2}\ge0.
\end{equation}
Lemma~\ref{lem:oc:hermite} proves \(f\le Q\) for \(-U\le a<U\).
Its remainder argument also covers \(a<-U\), since \(f'''\ge0\)
throughout the contact interval. Thus the stated majorant holds for
every \(a<U\), with \(B\ge0\).

If \(m\le0\), \eqref{eq:oc:fixed-contact} with \(h=1/2\) gives
\(\E e^{sW}\le(2/3)e^{-sU/2}+(1/3)e^{sU}\).
This centered two-point variable has interval length \(3U/2\),
so \eqref{eq:Q61:3.1} proves \eqref{eq:Q61:3.2} in this case.

If \(m\ge\allowbreak 0\), keep the variance
\(v=\allowbreak \E(W-m)^2=\allowbreak \E W^2-m^2\le\allowbreak  U^2/2-m^2\).
Then \(0\le\allowbreak  m\le\allowbreak  U/\sqrt2<\allowbreak U\). Put \(M=\allowbreak U-m>\allowbreak 0\) and
\(a=m-v/M\). Take the centered two-point variable
\(Y=\ell(Z-\pi)\) of \eqref{eq:oc:matched-contacts}, where
\(\ell=M+v/M\). Lemma~\ref{lem:oc:hermite} places \(a\) in
\([-U,U]\) and matches the first two moments, so
\begin{equation}\tag{Q61.3.4}\label{eq:Q61:3.4}
\E e^{sW}\le\E Q(W)=\E Q(m+Y)=e^{sm}\E e^{sY}.
\end{equation}
Its interval length obeys
\begin{equation}\tag{Q61.3.5}\label{eq:Q61:3.5}
M+\frac vM\le U-m+\frac{U^2/2-m^2}{U-m}
=\frac{3U^2/2-2Um}{U-m}\le\frac{3U}{2}.
\end{equation}
After multiplication by \(U-m>\allowbreak 0\), the last inequality is
\(-2Um\le\allowbreak -3Um/2\), valid for \(m\ge\allowbreak 0\).
Apply \eqref{eq:Q61:3.1} in \eqref{eq:Q61:3.4} to complete the proof.

\paragraph{Q61.4 Positive propagation and a uniform exponential payment.}
The scalar estimate on the whole interval \(|u|\le\allowbreak 1\) is
\begin{equation}\tag{Q61.4.1}\label{eq:Q61:4.1}
e^u\le1+u+\frac34u^2.
\end{equation}
For \(u\le\allowbreak 0\), the integral remainder coefficient is at most \(1/2\).
For \(0\le\allowbreak  u\le\allowbreak 1\), successive terms starting at degree two have
ratio at most \(1/3\), so their sum is at most
\((u^2/2)/(1-1/3)=\allowbreak 3u^2/4\).
Suppose a nonnegative message \(F_0\) has
\begin{equation}\tag{Q61.4.2}\label{eq:Q61:4.2}
F_0\le e^{R-\Lambda d},\qquad R=KP(K)g,\qquad
\|P(K)\|^2\le B,\qquad Bd\le1,
\end{equation}
where \(P\) is its finite arm polynomial. Set
\(V=\allowbreak \E R^2\), \(Z=\allowbreak KR\), \(U^2=\allowbreak Bd/3\).
The row, spectral and \(K^2\) row estimates give
\begin{equation}\tag{Q61.4.3}\label{eq:Q61:4.3}
\begin{gathered}
|R|^2\le Bd\le1,\quad V\le\theta Bd,\quad |Z|^2\le U^2,\\
\E Z^2\le\theta^2Bd=3\theta^2U^2\le U^2/2,
\end{gathered}
\end{equation}
using \(3\theta^2=\allowbreak 507/1024<\allowbreak 1/2\).
The kernel of \(T\) is nonnegative and at most two, so
\begin{equation}\tag{Q61.4.4}\label{eq:Q61:4.4}
Te^R\le q+\E R-Z+\tfrac32V=:B_0.
\end{equation}
Since \(B_0\ge\allowbreak  Te^R\ge\allowbreak 0\), AM--GM yields
\begin{equation}\tag{Q61.4.5}\label{eq:Q61:4.5}
0\le pTF_0\le e^{-\Lambda d}pB_0
\le e^{-\Lambda d}\left(1+\frac{\E R}{2}+\frac34V-\frac Z2\right)^2.
\end{equation}
The base equals \((p+B_0)/2\ge\allowbreak 0\). Applying \(x\le\allowbreak  e^{x-1}\),
taking the eighth power and integrating gives
\begin{equation}\tag{Q61.4.6}\label{eq:Q61:4.6}
\E[(pTF_0)^8]\le e^{-8\Lambda d+8\E R+12V}\E e^{-8Z}.
\end{equation}
Apply \eqref{eq:Q61:3.2} to the same \(W=\allowbreak -Z\), \(s=\allowbreak 8\):
\(\E e^{-8Z}\le\allowbreak  e^{8(-\E Z)_++18U^2}\).
Since \(12\theta+6=\allowbreak 87/8\), this gives
\begin{equation}\tag{Q61.4.7}\label{eq:Q61:4.7}
\boxed{\E[(pTF_0)^8]\le
\exp\left[-8\Lambda d+8\E R+\frac{87}{8}Bd+8(-\E Z)_+\right].}
\end{equation}

\paragraph{Q61.5.1 Finite pair envelopes.}
For \(a<\allowbreak b\), put \(R_{ab}=\allowbreak \sum_{j=\allowbreak a+1}^bK^j\one\),
\(\beta_1=\allowbreak \beta_2=\allowbreak 0\), and
\(\beta_b=\allowbreak \sum_{j=\allowbreak 1}^{\lfloor(b-1)/2\rfloor}(-1)^{j-1}\delta_j\).
Positive convolution gives the pair envelopes
\begin{equation}\tag{Q61.5.1}\label{eq:Q61:5.1}
h_af_b\le
\begin{cases}
e^{R_{ab}-\beta_b},&a=1,2,\\
e^{R_{ab}-(2d-\delta_2)},&a=3,\\
e^{R_{ab}-2d},&a\ge4.
\end{cases}
\end{equation}
For the first two short lengths, the pair-sum inequalities are
\(q+f_b\le\allowbreak 2-\beta_b+R_{1b}\) and
\(h_2+f_b\le\allowbreak 2-\beta_b+R_{2b}\).
For \(a=\allowbreak 3\), the sum is at most \(2+R_{3b}-(2d-\delta_2)\).
These are the finite short sums; the longer-pair estimate follows
also from the stronger bound \eqref{eq:L90:pair}.
AM--GM has a nonnegative base in each case.
For all \(b\ge\allowbreak 3\), \(\kappa d\le\allowbreak  d-\delta_2\le\allowbreak \beta_b\le\allowbreak  d\).
Removing \(E_0\) from \eqref{eq:Q61:0.1}, write the remaining message
as \(F_0\). Its exponent variable is
\begin{equation}\tag{Q61.5.2}\label{eq:Q61:5.2}
R=\sum_{i=1}^kR_{a_ib_i}=KP(K)g.
\end{equation}
Disjoint ordered intervals give a finite zero-one mask: the term
\(K^j\one\) corresponds to the coefficient of \(z^{j-2}\) in \(P\).

\paragraph{Q61.5.2 An arbitrary finite zero-one mask.}
At \(z=\allowbreak i\sqrt x\), \(0\le\allowbreak  x\le\allowbreak \theta\), the positive and negative
real subsums are each at most \(1/(1-x^2)\) in modulus; the imaginary
part is at most \(\sqrt x/(1-x^2)\). Hence
\begin{equation}\tag{Q61.5.3}\label{eq:Q61:5.3}
|P(i\sqrt x)|^2\le\frac{1+x}{(1-x^2)^2}\le\frac{54}{25}=:B_{\mathrm I}.
\end{equation}
Here \(1+\theta<\allowbreak 3/2\), \(\theta^2=\allowbreak 169/1024<\allowbreak 1/6\), and the
denominator exceeds \((5/6)^2\). This includes \(x=\allowbreak 0\).
The finite skew spectral theorem gives the corresponding operator norm.

\paragraph{Q61.5.3 The mask \(P_0=\allowbreak 1,\allowbreak P_1=\allowbreak P_2=\allowbreak 0\).}
Its positive real subsum is at most \(1/(1-x^2)\), the negative one
at most \(x^3/(1-x^2)\), and the imaginary part at most
\(x\sqrt x/(1-x^2)\). Thus
\begin{equation}\tag{Q61.5.4}\label{eq:Q61:5.4}
|P(i\sqrt x)|^2\le\frac{1+x^3}{(1-x^2)^2}
\le\frac{54}{35}<\frac{25}{16}=:B_{\mathrm {II}}.
\end{equation}
Indeed, \(\theta^2<\allowbreak 1/6\), \(\theta^3=\allowbreak 2197/32768<\allowbreak 1/14\), so the
ratio is below \((15/14)(36/25)=\allowbreak 54/35\); the last comparison is
\(864<\allowbreak 875\).

\paragraph{Q61.5.4 The mask \(P_0=\allowbreak P_2=\allowbreak 1,\allowbreak P_1=\allowbreak 0\).}
Here
\fitdisplay{
1-\frac{x}{1-x^2}\le\Re P\le1-x+\frac{x^2}{1-x^2}.
}{}
For \(0\le\allowbreak  x\le\allowbreak 1/2\), the lower bound is at least \(1/3\), and the
upper bound at most \(1-x/3\). The imaginary part is again at most
\(x\sqrt x/(1-x^2)\), so
\begin{equation}\tag{Q61.5.5}\label{eq:Q61:5.5}
|P|^2\le(1-x/3)^2+\frac{16}{9}x^3
\le(1-x/3)^2+\frac49x\le1-x/6\le1=:B_{\mathrm {III}}.
\end{equation}
The steps use \(x^2\le\allowbreak 1/4\) and \(x^2/9\le\allowbreak  x/18\).
These forced coefficients and gaps come from the ordered arm list.
The geometric sums bound finite subsums. All three constants are
at most \(54/25\), hence
\begin{equation}\tag{Q61.5.6}\label{eq:Q61:5.6}
Bd\le(54/25)(1/3)=18/25<1.
\end{equation}
Thus the exponential interval in Q61.4 holds on the whole domain.

\paragraph{Q61.6.1 The class \(b_1\ge\allowbreak 3\).}
All later \(a_i\ge\allowbreak 4\). Equation \eqref{eq:Q61:5.1} gives
\(F_0\le\allowbreak  e^{R-\Lambda d}\), \(\Lambda=\allowbreak \kappa+2(k-1)\).
The terms \(K^2\one,\allowbreak K^3\one\) both have coefficient one in \(R\).
Retain the \(-d\) contribution of the first and drop the other
negative mean terms:
\begin{equation}\tag{Q61.6.1}\label{eq:Q61:6.1}
\E R\le-d+\frac{\delta_2}{1-\theta^2}
\le-\frac{439}{855}d\le-d/2.
\end{equation}
Retain the \(+\delta_2\) contribution of the second to \(\E Z=\allowbreak \E KR\)
and drop the other positive terms:
\begin{equation}\tag{Q61.6.2}\label{eq:Q61:6.2}
\E Z\ge\delta_2-\frac{\delta_3}{1-\theta^2}
\ge\frac{439}{855}\delta_2\ge0.
\end{equation}
The coefficient used is \(\theta/(1-\theta^2)=\allowbreak 416/855\).
Insert \(B=\allowbreak B_{\mathrm I}=\allowbreak 54/25\) into \eqref{eq:Q61:4.7}:
\begin{equation}\tag{Q61.6.3}\label{eq:Q61:6.3}
\E[(pTF_0)^8]\le e^{-[16(k-2)+63/50]d}.
\end{equation}
The constant calculation is
\fitdisplay{
8\kappa+16+4-\frac{87}{8}\frac{54}{25}
=\frac{99}{4}-\frac{2349}{100}=\frac{63}{50}.
}{}
This holds for all \(b_1=\allowbreak 3\) or \(b_1\ge\allowbreak 4\), all later finite
lengths and \(0<\allowbreak d\le\allowbreak 1/3\).

\paragraph{Q61.6.2 The class \(b_1=\allowbreak 2,\allowbreak a_2\ge\allowbreak 4\).}
Every later short arm is at least four, so
\fitdisplay{
F_0\le e^{R-2(k-1)d},\qquad
R=K^2\one+\sum_{i=2}^k\sum_{j=a_i+1}^{b_i}K^j\one.
}{}
The latter sum starts at \(j\ge\allowbreak 5\), so \(P_0=\allowbreak 1,\allowbreak P_1=\allowbreak P_2=\allowbreak 0\),
and \(B=\allowbreak B_{\mathrm {II}}=\allowbreak 25/16\).
The first possible positive mean term has \(j=\allowbreak 8\); hence
\begin{equation}\tag{Q61.6.4}\label{eq:Q61:6.4}
\E R\le-d+\frac{\delta_4}{1-\theta^2}\le-11d/12.
\end{equation}
Here \(\theta^3/(1-\theta^2)=\allowbreak 2197/27360<\allowbreak 1/12\), by
\(26364<\allowbreak 27360\). The first possible negative term of \(\E Z\)
has \(j=\allowbreak 5\), giving
\begin{equation}\tag{Q61.6.5}\label{eq:Q61:6.5}
\E Z\ge-\frac{\delta_3}{1-\theta^2}
\ge-\frac{169}{855}d\ge-d/5,
\end{equation}
where \(845<\allowbreak 855\). Equation \eqref{eq:Q61:4.7} yields
\begin{equation}\tag{Q61.6.6}\label{eq:Q61:6.6}
\E[(pTF_0)^8]\le e^{-[16(k-2)+9103/1920]d}.
\end{equation}
The exact constant is
\fitdisplay{
16+\frac{22}{3}-\frac85-\frac{87}{8}\frac{25}{16}
=\frac{9103}{1920}>4>\frac{63}{50}.
}{}
The \(8/5\) term retains the possible negative mean of \(Z\).

\paragraph{Q61.6.3 The class \(b_1=\allowbreak 2,\allowbreak a_2=\allowbreak 3\).}
The direction remains
\(qf_2h_3f_{b_2}\prod_{i\ge\allowbreak 3}h_{a_i}f_{b_i}\).
The second pair has loss \(2d-\delta_2\), and each later pair loss
\(2d\), so \(\Lambda=\allowbreak 2(k-1)-\theta\).
The terms \(K^2\one,\allowbreak K^4\one\) are present, and \(K^3\one\) is absent.
Thus \(P_0=\allowbreak P_2=\allowbreak 1,\allowbreak P_1=\allowbreak 0\), and \(B=\allowbreak B_{\mathrm {III}}=\allowbreak 1\).
The same first-sign estimates give
\begin{equation}\tag{Q61.6.7}\label{eq:Q61:6.7}
\E R\le-d/2,\qquad\E Z\ge-d/5.
\end{equation}
Consequently
\begin{equation}\tag{Q61.6.8}\label{eq:Q61:6.8}
\E[(pTF_0)^8]\le e^{-[16(k-2)+171/40]d},
\end{equation}
since
\fitdisplay{
16-8\theta+4-\frac85-\frac{87}{8}
=\frac{171}{40}>4>\frac{63}{50}.
}{}
This includes \(k=\allowbreak 2\) and every \((1,\allowbreak 2,\allowbreak 3,\allowbreak c)\) in this direction.

\paragraph{Q61.6.4 Completion of the three classes.}
For \(k\ge\allowbreak 2\), either \(b_1\ge\allowbreak 3\), or \(b_1=\allowbreak 2\) with \(a_2=\allowbreak 3\)
or \(a_2\ge\allowbreak 4\). Each class has the stronger bound
\(\E[(pTF_0)^8]\le\allowbreak  e^{-[16(k-2)+63/50]d}\), and hence
\(\E[(pTF_0)^8]\le\allowbreak  e^{-63d/50}\).
Using \(e^x\ge\allowbreak 1+x\) and \(d\le\allowbreak 1/3\),
\begin{equation}\tag{Q61.6.9}\label{eq:Q61:6.9}
e^{-63d/50}\le\frac1{1+63d/50}
=1-\frac{63d}{50+63d}\le1-\frac{63}{71}d\le1-d/2.
\end{equation}
Restore \(E_0\le\allowbreak 1\) by
\(0\le\allowbreak  pT(E_0F_0)\le\allowbreak  pTF_0\), proving \eqref{eq:Q61:0.2}.

\paragraph{Q61.7 Decorations, extreme entries and full reversal.}
Genuine equal pairs at both terminal and receiving root multiply
nonnegative messages by factors at most one. Positivity restores
every finite multiplicity and pair length. Entries \(|K(x,\allowbreak y)|=\allowbreak 1\)
are allowed; the proof divides by neither \(p\) nor \(q\).
The one-atom space has \(d=\allowbreak 0\); a zero message is also compatible
with AM--GM and the exponential upper bounds.
If \(\overline F\) is obtained by reversing every edge of \(F\),
replace \(K\) by \(-K\) in the whole formula. Then \(p\) becomes
\(q\), \(T\) becomes \(A\), and \(d,\allowbreak \delta_j\) stay unchanged. Thus
\begin{equation}\tag{Q61.7.1}\label{eq:Q61:7.1}
\E[(qA\overline F)^8]\le e^{-63d/50}\le1-d/2.
\end{equation}
All terminal, connector, receiving-unit and decoration edges reverse
together. The ordered terminal direction is also that of
Appendix~\ref{loc:S29}, which independently proves
\begin{equation}\tag{Q61.7.2}\label{eq:Q61:7.2}
\boxed{\E[p(T^nF)^4]\le1\qquad(n\ge1).}
\end{equation}
